\subsection{局部道路连通空间}

定义 4. --- 若拓扑空间的每一点都有由道路连通邻域组成的基本邻域系，则称它是局部道路连通的。

命题 7. --- 局部道路连通拓扑空间是局部连通的。其连通分支与道路连通分支重合。特别地，若局部道路连通拓扑空间连通，则它道路连通。

第一个断言由命题 4 得出。我们来证明第二个断言。设 X 为局部道路连通拓扑空间，C 为 X 的一个道路连通分支。C 的每一点都有道路连通的邻域，它包含于 C；因此 C 是 X 的开子集。由于诸道路连通分支构成 X 的一个分割，这样的分支也是闭的。由于它连通（III, p. 258, 命题 4），它是一个连通分支（TG, I, p. 83）。第三个断言由第二个得出。

于是注意，说一个拓扑空间连通且局部道路连通并无歧义。

推论 1. --- 局部道路连通拓扑空间的每个连通开子集是道路连通的。

推论 2. --- 设 B 为局部道路连通拓扑空间，E 为平展 B-空间。空间 E 是局部道路连通的。若它连通，则它道路连通。

第一个断言由定义 4 及平展态射的定义（I, p. 28, 定义 2）立即得出。第二个由命题 7 从它推出。

为使拓扑空间 X 局部道路连通，必须而且只需 X 的开子集的每个道路连通分支都是 X 的开子集（参见 I, p. 85, 命题 11 的证明）。若空间 X 局部道路连通，则 X 的每一点都有由开的道路连通邻域组成的基本邻域系。

命题 8. --- 局部道路连通空间的每个商空间都是局部道路连通的。

设 X 为局部道路连通空间，R 为 X 上的等价关系，\(\varphi\colon\mathrm{X}\to\mathrm{X}/\mathrm{R}\) 为典范映射。只需证明 X/R 的开子集的道路连通分支都是开的。设 A 为 X/R 的开子集，C 为 A 的一个道路连通分支。设 \(x\in\bar{\varphi}^{1}(\mathrm{C})\)，K 为 x 在 \(\bar{\varphi}^{1}(\mathrm{A})\) 中的道路连通分支。集合 \(\varphi(\mathrm{K})\) 是道路连通的（III, p. 258, 命题 3），包含于 A 且包含 \(\varphi(x)\)；于是有 \(\varphi(\mathrm{K})\subset\mathrm{C}\)，从而 \(\mathrm{K}\subset\bar{\varphi}^{1}(\mathrm{C})\)。这证明 \(\bar{\varphi}^{1}(\mathrm{C})\) 是 \(\bar{\varphi}^{1}(\mathrm{A})\) 的道路连通分支的并。由于 X 局部道路连通且 \(\bar{\varphi}^{1}(\mathrm{A})\) 在 X 中开，\(\bar{\varphi}^{1}(\mathrm{C})\) 在 X 中开。因此 C 在 X/R 中开。命题得证。

命题 9. --- 一族局部道路连通空间（除有限个外都连通）的积是局部道路连通的。

设 \((\mathrm{X}_{j})_{j\in\mathrm{J}}\) 为一族局部道路连通空间，除有限个外都连通。设 \(x=(x_{j})_{j\in\mathrm{J}}\) 为积空间 \(X=\prod_{j\in J}X_{j}\) 的一点。按假设，x 在 X 中的基本邻域系由形如 \(V=\prod_{j\in J}V_{j}\) 的集合组成，其中对每个 \(j\in J\)，\(V_{j}\) 是 \(x_{j}\) 在 \(X_{j}\) 中的道路连通邻域，且 \(V_{j}=X_{j}\)（除有限多个指标 \(j\in J\) 外）（TG, I, p. 24）。这些集合都是道路连通的（III, p. 260, 命题 6），故 X 局部道路连通。

推论. --- 数值空间的每个开子集都是局部道路连通的。

R 的每一点都有由区间组成的邻域基；因此 R 是局部道路连通的。由命题 9，数值空间 \(R^{n}\)（对每个整数 \(n \geqslant 1\)）都是局部道路连通的。此外，局部道路连通拓扑空间的开子集仍是局部道路连通的，推论得证。

例. --- 设 G 为 \(\mathbf{GL}(n,\mathbf{R})\) 中由行列式严格正的 n 阶方阵组成的子群。群 G 是连通的且局部道路连通的。

由 A, III, p. 104, 命题 17，群 \(\mathbf{SL}(n,\mathbf{R})\) 由元素 \(B_{ij}(\lambda)\)（\(1\leqslant i,j\leqslant n\)，\(i\neq j\)，\(\lambda\in\mathbf{R}\)）生成。映射 \(\lambda\mapsto B_{ij}(\lambda)\)（从 \(\mathbf{R}\) 到 \(\mathbf{GL}(n,\mathbf{R})\) 中）是连续的，它们的像是 \(\mathbf{SL}(n,\mathbf{R})\) 的连通子集；由于它们都包含单位矩阵 \(I_n\)，它们的并是连通的（TG, I, p. 81, 命题 2）。因此群 \(\mathbf{SL}(n,\mathbf{R})\) 是连通的（TG, III, p. 8, 命题 7）。设 A 为 G 中由形如 diag(1, \(\ldots\),1, \(\lambda\))（\(\lambda\in\mathbf{R}_{+}^{*}\)）的矩阵组成的子群；它是连通的，且有 \(\mathbf{GL}(n,\mathbf{R})=\mathbf{A}\cdot\mathbf{SL}(n,\mathbf{R})\)。由此推出群 G 是连通的。由于它是 \(\mathbf{R}_{+}^{*}\) 在从 \(\mathbf{M}_{n}(\mathbf{R})\) 到 \(\mathbf{R}\) 中的行列式映射下的原像，它是 \(\mathbf{M}_{n}(\mathbf{R})\) 的开子集；因此它是局部道路连通的（前述推论），也是道路连通的（III, p. 261, 推论 1）。

命题 10. --- 设 X 为拓扑空间。映射 \((o, e)\)（从 \(\Lambda(\mathrm{X}) = \mathcal{C}_{c}(\mathbf{I}; \mathrm{X})\) 到 \(X \times X\) 中，它把 X 中的道路 c 对应到对 \((c(0), c(1))\)）是连续的。若空间 X 局部道路连通，则该映射是开的。

我们已经知道 \(\Lambda(\mathrm{X})\) 到 X 中的起点映射 o 与终点映射 e 是连续的（III, p. 257），第一个断言由此得出。

引理. --- 设 \(c\colon \mathbf{I}\to \mathrm{X}\) 为道路，W 为 \(c\) 在空间 \(\mathcal{C}_c(\mathbf{I};\mathrm{X})\) 中的邻域。存在实数 \(\varepsilon \in ]0,1 / 2]\)、邻域 \(\mathrm{V}_0\)（含 \(c(0)\)）与邻域 \(\mathrm{V}_1\)（在 X 中、含 \(c(1)\)），具有下列性质：我们有 \(c([0,\varepsilon])\subset \mathrm{V}_0\)，\(c([1 - \varepsilon,1])\subset \mathrm{V}_1\)，且每条道路 \(c^{\prime}\colon \mathbf{I}\rightarrow \mathrm{X}\)，若它满足

\[
\left\{ \begin{array}{l l} c ^ {\prime} (t) \in \mathrm{V} _ {0} & \text { for } 0 \leqslant t \leqslant \varepsilon , \\ c ^ {\prime} (t) = c (t) & \text { for } \varepsilon \leqslant t \leqslant 1 - \varepsilon , \\ c ^ {\prime} (t) \in \mathrm{V} _ {1} & \text { for } 1 - \varepsilon \leqslant t \leqslant 1, \end{array} \right.\tag{2}
\]

，便属于 W。

由紧收敛拓扑的定义（TG, X, p. 26, 定义 1），存在有限集 J、X 的开集族 \((\mathrm{U}_{j})_{j\in\mathrm{J}}\) 与 I 的紧子集族 \((\mathrm{K}_{j})_{j\in\mathrm{J}}\)，使得集合 \(W'\)（其中的道路 \(c'\) 对每个指标 j 满足 \(c'(\mathrm{K}_{j})\subset\mathrm{U}_{j}\)）是包含于 W 中的 c 的邻域。然后用 \(A_{0}\)（相应地，\(A_{1}\)）表示满足 \(0\in K_{j}\)（相应地，\(1\in K_{j}\)）的指标 j 的集合；令 \(V_{0}=\bigcap_{j\in A_{0}}U_{j}\)，\(V_{1}=\bigcap_{j\in A_{1}}U_{j}\)。

由于映射 \(c\) 连续，存在实数 \(\varepsilon \in [0,1/2]\)，使得 \(c([0,\varepsilon]) \subset \mathrm{V}_0\)，\(c([1-\varepsilon,1]) \subset \mathrm{V}_1\)，且 \([0,\varepsilon] \cap \mathrm{K}_j = \emptyset\)（对所有 \(j \notin \mathrm{A}_0\)），\([1-\varepsilon,1] \cap \mathrm{K}_j = \emptyset\)（对所有 \(j \notin \mathrm{A}_1\)）。设 \(c'\) 为满足条件 (2) 的道路。我们来证明 \(c' \in \mathrm{W}'\)。设 \(j \in \mathrm{J}\)，\(t \in \mathrm{K}_j\)。若 \(\varepsilon \leqslant t \leqslant 1-\varepsilon\)，则 \(c'(t) = c(t)\) 属于 \(\mathrm{U}_j\)。若 \(0 \leqslant t \leqslant \varepsilon\)，则 \(c'(t) \in \mathrm{V}_0\)；由 \(\varepsilon\) 的取法有 \(j \in \mathrm{A}_0\)，因此 \(c'(t) \in \mathrm{U}_j\)。类似地，若 \(1-\varepsilon \leqslant t \leqslant 1\)，则有 \(j \in \mathrm{A}_1\) 且 \(c'(t) \in \mathrm{V}_1 \subset \mathrm{U}_j\)。于是 \(c'(\mathrm{K}_j) \subset \mathrm{U}_j\)，从而 \(c'\) 属于 \(\mathrm{W}'\)，故属于 W。

现在证明命题 10 的第二个断言。设空间 X 局部道路连通。设 \(c\) 为 X 中的道路，W 为 \(c\) 在 \(\mathcal{C}_c(\mathbf{I};\mathbf{X})\) 中的邻域。取引理中的 \(\varepsilon\)、\(\mathrm{V}_0\) 与 \(\mathrm{V}_1\)。设 \(\mathrm{T}_0\) 与 \(\mathrm{T}_1\) 为 \(c(0)\) 与 \(c(1)\) 的道路连通邻域，分别包含于 \(\mathrm{V}_0\) 与 \(\mathrm{V}_1\) 中。存在实数 \(\theta\)，满足 \(0 < \theta < \varepsilon\)，且 \(c([0,\theta]) \subset \mathrm{T}_0\)，\(c([1 - \theta,1]) \subset \mathrm{T}_1\)。设 \(x_0 \in \mathrm{T}_0\)，\(x_1 \in \mathrm{T}_1\)；设 \(c_0\) 为以 \(x_0\) 为起点、\(c(\theta)\) 为终点的道路（在 \(\mathrm{T}_0\) 中），\(c_1\) 为以 \(x_1\) 为起点、\(c(1 - \theta)\) 为终点的道路

（在 \(\mathrm{T}_1\) 中）。令

\[
c ^ {\prime} (t) = \left\{ \begin{array}{l l} c _ {0} (t / \theta) & \text { for } 0 \leqslant t \leqslant \theta , \\ c (t) & \text { for } \theta \leqslant t \leqslant 1 - \theta , \\ c _ {1} ((1 - t) / \theta) & \text { for } 1 - \theta \leqslant t \leqslant 1. \end{array} \right.
\]

这样便定义了一条道路 \(c'\)，它连接 \(x_{0}\) 与 \(x_{1}\) 并满足条件 (2)。这证明 W 在映射 \((o, e)\) 下在 X × X 中的像包含邻域 \(T_{0} \times T_{1}\)（\((c(0), c(1))\) 的），命题得证。

推论. --- 设 X 为局部道路连通拓扑空间，x 为 X 的一点。映射 \(c \mapsto c(1)\)（从 \(\Lambda_{x}(\mathrm{X})\) 到 X 中）是开的。

由本命题，映射 \(\varphi\colon\Lambda(\mathrm{X})\to\mathrm{X}\times\mathrm{X}\)（由 \(\varphi(c)=(c(0),c(1))\) 定义）是开的，且有 \(\Lambda_{x}(\mathrm{X})=\overline{\varphi}^{1}(\{x\}\times\mathrm{X})\)。因此，映射 \(\Lambda_{x}(\mathrm{X})\to\{x\}\times\mathrm{X}\)（由 \(\varphi\) 诱导）是开的（TG, I, p. 30, 命题 2），它与第二投影 \(\mathrm{pr}_{2}\) 的复合亦然。

